\clearpage
\section{Recognizing Sound Shapes}\label{sec:shapes}

A spectrogram turns musical events into shapes. Read position as frequency,
horizontal extent as captured time, and color as magnitude. Start with
Average and Smooth off so their blending does not disguise the pattern.

\begin{figure}[!htp]
\centering
% Abstract time-frequency signatures; no app data or simulated UI.
\begin{tikzpicture}[x=1cm,y=1cm,font=\small]
  \foreach \x/\y/\title in {0/3.4/Steady harmonics,7/3.4/Pitch glide,0/0/Attack and decay,7/0/Filtered noise} {
    \begin{scope}[shift={(\x,\y)}]
      \fill[panelPaper] (0,0) rectangle (5.5,2.1);
      \node[anchor=west,font=\small\bfseries,text=spectreAccent] at (0,2.5) {\title};
      \draw[->,panelMuted] (0,0)--(5.7,0) node[right,font=\scriptsize] {$t$};
      \draw[->,panelMuted] (0,0)--(0,2.2) node[above,font=\scriptsize] {$f$};
    \end{scope}
  }
  \foreach \y/\a in {.45/90,.95/65,1.45/40} {
    \draw[spectreAccent!\a,line width=2pt] (.3,3.4+\y)--(5.1,3.4+\y);
  }
  \foreach \y/\a in {.2/90,.65/60,1.1/35} {
    \draw[spectreAccent!\a,line width=2pt] (7.3,3.4+\y) .. controls
      (8.5,3.4+\y) and (9,3.8+\y) .. (12.1,4.3+\y);
  }
  \fill[spectreAccent!70] (.6,.15) rectangle (.85,1.95);
  \foreach \i in {0,...,15} {
    \pgfmathtruncatemacro{\shade}{85-\i*5}
    \draw[spectreAccent!\shade,line width=2.5pt] (.85+\i*.25,.6)--(1.1+\i*.25,.6);
    \draw[spectreAccent!\shade,line width=1.5pt] (.85+\i*.19,1.2)--(1.04+\i*.19,1.2);
  }
  \fill[spectreAccent!12] (7.3,.15)--(7.3,.6).. controls (9,.6) and (9,1.8)..(12.1,1.8)--(12.1,.15)--cycle;
  \draw[spectreAccent,line width=1.6pt] (7.3,.6)..controls (9,.6) and (9,1.8)..(12.1,1.8);
  \foreach \i in {0,...,22} {
    \pgfmathsetmacro{\xx}{7.35+\i*.21}
    \draw[spectreAccent!25,line width=.6pt] (\xx,.22)--(\xx,.38);
  }
\end{tikzpicture}

\caption{Schematic sound signatures, shown in chronological order within
one sweep. Marks indicate energy; their shade is illustrative rather than a
calibrated palette. Spectre's scan line wraps when it reaches the right edge.}
\end{figure}

\begin{description}[style=nextline,leftmargin=0pt,labelwidth=0pt,labelsep=0pt]
  \item[Steady tone] A sine produces one horizontal band. A harmonic sound
  adds bands near integer multiples of its fundamental. With Linear Freq
  Scale, equally spaced harmonics are equally spaced on the frequency axis.
  \item[Pitch movement] A glide bends the bands upward or downward.
  Vibrato makes them undulate. Harmonics move together as the fundamental
  changes; frequency smoothing can blend neighboring bands.
  \item[Attack and decay] A short broadband attack produces a vertical
  mark. Resonances can continue as horizontal tails. A pitched percussion
  sound may combine a vertical attack with a falling pitch curve.
  \item[Sustained noise] Energy spreads over a frequency region instead
  of forming stable harmonic lines. A moving filter changes that region's
  boundary and may leave a bright resonant band near its cutoff.
\end{description}

\paragraph{Use shape as evidence}
A thick band can be a window's broad peak, modulation, several nearby tones,
or smoothing. It does not establish that the source is noisy. An evenly
colored region can simply be above Ceil or below Floor. Change one viewing
control at a time, then compare the sound again with minimal smoothing.
The next walkthroughs isolate these causes.

\clearpage
\section{Following An Envelope}\label{sec:envelope}

Follow how a sound develops rather than only how bright its final spectrum
looks. Use an oscillator or noise source, a VCA, and a repeating envelope.

\paragraph{Starting point}
Load \textbf{DecibelInspection}. It sets unity input gain, Flattop, no
smoothing, zero Slope, and a $-90$ to $+12\,\mathrm{dB}$ color range.
Change Window to Hann. Use Linear Freq Scale and $0$--$5\,\mathrm{kHz}$
for a pitched source, or a wider range for noise. Keep Rack at
$48\,\mathrm{kHz}$ for the timing examples here.

\begin{enumerate}
  \item Patch the sound through the VCA. Send the envelope to the VCA's
  level input, and branch the VCA output to Spectre and your listening path.
  Trigger a short attack and a decay of roughly half a second once a second.
  \item Watch several repetitions. A sine leaves one band whose color
  brightens at each attack and fades during decay. Noise leaves a broad
  vertical event. Your envelope shape controls how those colors fade.
  \item Set Average to $300\,\mathrm{ms}$ and observe new events. Their
  level changes should look more gradual. Previously recorded columns are
  unchanged; Average affects newly analyzed spectra. Restore Average to zero.
  \item Send an envelope to oscillator pitch or to a filter's cutoff.
  Compare movement in frequency with fading amplitude. Freeze after a few
  events and inspect their shapes; resume to try a different envelope.
\end{enumerate}

\paragraph{What time detail survives?}
At $48\,\mathrm{kHz}$ each column summarizes about $42.7\,\mathrm{ms}$
of audio, and successive windows overlap by half. A very short click can
therefore affect several columns. Its drawn width is not its exact duration.
Averaging adds further persistence. Use an oscilloscope for individual sample
edges or precise attack timing, and Spectre for the changing frequency content.

\paragraph{Try a fair comparison}
Keep gain, Floor, Ceil, window, and sample rate fixed while changing the
source. A darker decay after raising Floor does not show that the VCA became
quieter: it shows that you hid more low-level detail. Likewise, changing
Slope can make upper partials appear to sustain longer by making them more
visible. Use zero Slope while comparing decay across frequencies.

\clearpage
\section{Finding Quiet Detail}\label{sec:detail}

Separate an acquisition change from a viewing change by recoloring the same
frozen history. Use a sound with both strong and weak partials or a long decay.

\paragraph{Starting point}
Load \textbf{DecibelInspection} and capture a few seconds of the sound.
Keep input gain at $0\,\mathrm{dB}$ and Slope at zero. Then press Run to freeze.

\begin{enumerate}
  \item Hover over a strong band and a quieter region. In Decibels mode,
  Raw reports the stored bin below the cursor frequency; Color reports the
  weighted, interpolated image sample. Neither readout is clamped to the
  palette's endpoints.
  \item Raise Floor from $-90$ to $-60\,\mathrm{dB}$ using its right-click
  numeric entry. Weak detail now shares the bottom color. This changes
  contrast; it does not remove noise or rewrite the stored spectra.
  \item Lower Floor again and lower Ceil toward the quiet detail. More of
  the palette is allocated to that range, while strong peaks can saturate at
  the top color. Raise Ceil again to distinguish the strong values.
  \item Try Cividis or Gray to compare color and lightness cues. Switch to
  Linear and back to Decibels. Each scale restores its own range; the stored
  history is the same throughout the experiment.
\end{enumerate}

\paragraph{A worked contrast example}
With Floor $-90$ and Ceil $0\,\mathrm{dB}$, a $-60\,\mathrm{dB}$ sample
sits one third of the way through the palette. With Floor $-70$ and Ceil
$-50\,\mathrm{dB}$, that same sample sits halfway through it. Values above
$-50\,\mathrm{dB}$ all reach the top color in the second view. Palette
position is numeric; equal steps need not look equally bright in every map.

\begin{center}
\renewcommand{\arraystretch}{1.2}
\begin{tabularx}{\textwidth}{@{}p{.35\textwidth}X@{}}
\toprule
\textbf{Change while frozen} & \textbf{Effect} \\
\midrule
Palette, range, Slope & Recolors retained history. \\
Freq Scale, LO/HI Freq & Remaps or crops the frequency view. \\
Window, Average, Smooth & Takes effect on new analysis after resuming. \\
Input gain, AC coupling & Affects samples captured after resuming. \\
\bottomrule
\end{tabularx}
\end{center}

To compare analysis settings, resume and allow fresh history to replace the
old columns. Save an external image to preserve a result: a saved Rack patch
retains settings and Run state, but not the captured history.

\clearpage
\section{Time And Measurement Limits}\label{sec:time}

\paragraph{Read the sweep as a circular timeline}
After the first full sweep, the oldest column is just ahead of the scan line
and the newest is just behind it. To follow events from old to new, start at
the scan line, read rightward to the edge, then continue from the left edge
back to the scan line. Events on opposite sides of the scan line are almost
one history length apart, despite being neighbors on screen.

\begin{center}
\renewcommand{\arraystretch}{1.25}
\begin{tabular}{@{}rrrr@{}}
\toprule
\textbf{Rack rate} & \textbf{Bin spacing} & \textbf{Column interval} & \textbf{Full sweep} \\
\midrule
$44.1\,\mathrm{kHz}$ & $21.53\,\mathrm{Hz}$ & $23.22\,\mathrm{ms}$ & $11.89\,\mathrm{s}$ \\
$48\,\mathrm{kHz}$ & $23.44\,\mathrm{Hz}$ & $21.33\,\mathrm{ms}$ & $10.92\,\mathrm{s}$ \\
$96\,\mathrm{kHz}$ & $46.88\,\mathrm{Hz}$ & $10.67\,\mathrm{ms}$ & $5.46\,\mathrm{s}$ \\
\bottomrule
\end{tabular}
\end{center}

The analysis window lasts twice the column interval. At $48\,\mathrm{kHz}$,
one quarter of the history width represents about $2.73\,\mathrm{s}$ of
captured time. This is an estimate, not a cursor time measurement. Pauses are
excluded, and resuming can finish a window that began before the pause; the
first resumed columns need not represent uninterrupted real-world time.
After a sample-rate change, wait a full sweep before using one time scale
for the entire image.

\paragraph{Know what a level means}
For a steady isolated sine, unity gain, zero Slope, no smoothing, AC coupling
off, and a bin-centered frequency, $5\,\mathrm{V}$ peak is approximately
$0\,\mathrm{dB}$ in Decibels mode. Half that amplitude is about
$-6.02\,\mathrm{dB}$; one tenth is $-20\,\mathrm{dB}$. At
$48\,\mathrm{kHz}$, $750\,\mathrm{Hz}$ is exactly bin 32 of the fixed
2048-sample FFT and provides a convenient reference tone. Inspect Raw near
that bin; between bins, the Color readout can differ because it interpolates.

\paragraph{Know what a level cannot tell you}
The largest spectral bin is not total RMS, perceived loudness, or noise power
per hertz. A finite window spreads a tone over neighboring bins. Quiet bands
near a strong tone can be window leakage; changing window requires new
capture to test that possibility. Spectre's fixed length cannot resolve every
pair of nearby bass tones: cropping magnifies existing detail. Use Fourier's
longer frames when frequency separation matters more than history.

\paragraph{Keep sources separate when needed}
A polyphonic input sums voices before analysis. Identical in-phase voices
reinforce; opposite-polarity copies can cancel. The image contains no phase
measurement. Split voices before Spectre to inspect one voice, or use separate
Fourier inputs to compare several signals at once.
